\chapter{Common Backend Errors}

\begin{notebox}[NOTE]
Same reference format as Chapter 41: \textbf{Cause}, \textbf{Identify},
\textbf{Solution} for each error, in the order they tend to appear while
building the backend.
\end{notebox}

\subsubsection*{"Cannot find module 'X'"}
\noindent\begin{tabularx}{\textwidth}{@{}>{\bfseries}p{2.3cm}X@{}}
\toprule
Cause & The package isn't installed, the import path is misspelled, or a
relative import is missing its file extension when using ESM. \\
Identify & Node throws \texttt{Error: Cannot find module} immediately on
startup, with the path it tried to resolve. \\
Solution & Run \texttt{npm install <package>} for third-party modules, or
fix the relative path for local files. \\
\bottomrule
\end{tabularx}

\subsubsection*{Port Already in Use (EADDRINUSE)}
\noindent\begin{tabularx}{\textwidth}{@{}>{\bfseries}p{2.3cm}X@{}}
\toprule
Cause & Another process (often a previous, still-running \texttt{nodemon}
instance) is already bound to the same port. \\
Identify & \texttt{Error: listen EADDRINUSE: address already in use
:::5000} on startup. \\
Solution & Kill the process on that port or change \texttt{PORT} in
\texttt{.env}, then restart. \\
\bottomrule
\end{tabularx}

\subsubsection*{MongoDB Connection Failed}
\noindent\begin{tabularx}{\textwidth}{@{}>{\bfseries}p{2.3cm}X@{}}
\toprule
Cause & Wrong connection string, wrong credentials, or the current IP
isn't on the Atlas Network Access allowlist. \\
Identify & \texttt{MongooseServerSelectionError} or
\texttt{MongoNetworkError} shortly after \texttt{mongoose.connect()}. \\
Solution & Re-check \texttt{MONGO\_URI}, confirm the Atlas user's
password, and add your IP (or \texttt{0.0.0.0/0} for testing) to Network
Access. \\
\bottomrule
\end{tabularx}

\subsubsection*{req.body Undefined}
\noindent\begin{tabularx}{\textwidth}{@{}>{\bfseries}p{2.3cm}X@{}}
\toprule
Cause & \texttt{express.json()} middleware isn't registered, or is
registered after the routes that need it. \\
Identify & \texttt{req.body} is \texttt{undefined} even though the
frontend sent a JSON payload (visible in the Network tab request body). \\
Solution & Add \texttt{app.use(express.json())} in \texttt{app.js}
\emph{before} \texttt{app.use('/api', router)}. \\
\bottomrule
\end{tabularx}

\subsubsection*{req.user Undefined}
\noindent\begin{tabularx}{\textwidth}{@{}>{\bfseries}p{2.3cm}X@{}}
\toprule
Cause & The \texttt{protect} middleware wasn't attached to the route, or
was attached after the controller in the middleware chain. \\
Identify & Controller throws when reading \texttt{req.user.\_id} even
though the request included a valid cookie/token. \\
Solution & Add \texttt{protect} as the route's first argument:
\texttt{router.get('/me', protect, getMe)} (Chapter 20). \\
\bottomrule
\end{tabularx}

\subsubsection*{JWT Invalid}
\noindent\begin{tabularx}{\textwidth}{@{}>{\bfseries}p{2.3cm}X@{}}
\toprule
Cause & \texttt{JWT\_SECRET} differs between signing and verifying (e.g.
different \texttt{.env} on a restarted server), or the token was
tampered with. \\
Identify & \texttt{jwt.verify()} throws
\texttt{JsonWebTokenError: invalid signature}. \\
Solution & Ensure the same \texttt{JWT\_SECRET} is used everywhere and
was not rotated without invalidating old sessions. \\
\bottomrule
\end{tabularx}

\subsubsection*{JWT Expired}
\noindent\begin{tabularx}{\textwidth}{@{}>{\bfseries}p{2.3cm}X@{}}
\toprule
Cause & The token's \texttt{exp} claim has passed --- normal behavior for
a token older than its \texttt{expiresIn} setting. \\
Identify & \texttt{jwt.verify()} throws \texttt{TokenExpiredError}. \\
Solution & Respond with \texttt{401} so the frontend can redirect to
login, or implement a refresh-token flow if silent renewal is needed. \\
\bottomrule
\end{tabularx}

\subsubsection*{Mongoose ValidationError}
\noindent\begin{tabularx}{\textwidth}{@{}>{\bfseries}p{2.3cm}X@{}}
\toprule
Cause & A document failed a schema rule --- missing \texttt{required}
field, value outside \texttt{enum}, wrong type, etc. \\
Identify & Error object has \texttt{name: 'ValidationError'} and an
\texttt{errors} map keyed by field name. \\
Solution & Return \texttt{400} with the specific field messages; fix the
payload or relax/correct the schema rule. \\
\bottomrule
\end{tabularx}

\subsubsection*{Duplicate Key Error (E11000)}
\noindent\begin{tabularx}{\textwidth}{@{}>{\bfseries}p{2.3cm}X@{}}
\toprule
Cause & Insert/update violates a \texttt{unique} index --- most commonly
a duplicate \texttt{email} on the User model. \\
Identify & Error message contains \texttt{E11000 duplicate key error}
and the offending field/value. \\
Solution & Catch the error code \texttt{11000} in the controller and
return a \texttt{400} like "Email already registered." \\
\bottomrule
\end{tabularx}

\subsubsection*{500 Internal Server Error}
\noindent\begin{tabularx}{\textwidth}{@{}>{\bfseries}p{2.3cm}X@{}}
\toprule
Cause & Generic catch-all --- an unhandled exception somewhere in the
request lifecycle (bad destructuring, thrown error with no specific
handling, etc.). \\
Identify & Frontend just sees "500"; the \emph{real} cause is only visible
in the backend terminal's stack trace. \\
Solution & Read the stack trace top to bottom, find the first line
pointing at your own code (not \texttt{node\_modules}), and fix that. \\
\bottomrule
\end{tabularx}

\subsubsection*{Route Not Found / 404}
\noindent\begin{tabularx}{\textwidth}{@{}>{\bfseries}p{2.3cm}X@{}}
\toprule
Cause & Typo in the route path, wrong HTTP method, or a more general
route defined earlier in the file is intercepting the request. \\
Identify & Falls through to Express's default 404 handler (or your
catch-all) instead of the intended controller. \\
Solution & Double-check the method + path match exactly, and that
specific routes (\texttt{/tasks/search}) are declared before parameterized
ones (\texttt{/tasks/:id}). \\
\bottomrule
\end{tabularx}

\subsubsection*{CORS Error (Backend Misconfiguration)}
\noindent\begin{tabularx}{\textwidth}{@{}>{\bfseries}p{2.3cm}X@{}}
\toprule
Cause & \texttt{cors()} is missing, uses the wildcard \texttt{origin: '*'}
alongside \texttt{credentials: true} (invalid combination), or lists the
wrong frontend domain. \\
Identify & Preflight \texttt{OPTIONS} request in the Network tab returns
without the expected \texttt{Access-Control-Allow-*} headers. \\
Solution & Set an explicit origin and \texttt{credentials: true} together;
never combine \texttt{'*'} with credentials (Chapter 26). \\
\bottomrule
\end{tabularx}
